\section{Lecture 1 (September 2nd)}
\begin{rmk}
The main subject of this course is Galois theory. The motivation of this theory is to find a general solution of polynomial equations. For example, the general solution of the quadratic equation
\[ax^2+bx+c\;=\;0\]
where $a\ne 0$, $b$ and $c$ are real numbers is given by
\[x\;=\;\dfrac{-b\pm\sqrt{b^2-4ac}}{2a}\]
To be more precise, we want to express a general solution of the equation
\[a_{n}x^{n}+a_{n-1}x^{n-1}+\ldots +a_{0}=0\]
where $a_{i}\in {\bm Q}$ and $a_{n}\ne 0$ using 4 basic operations (namely, $\pm$, $\times $, and $\div$), and $n$-th roots. Today, we know that, through Galois theory, this is not possible for $n\geq 5$. 

\bigskip

In essence, Galois theory connects the theory of fields and the theory of groups. Consider the following example. The polynomial
\[x^4+x^3+x^2+x+1=0\]
We know how to solve this,
\[(x-1)(x^4+x^3+x^2+x+1)=x^5-1=0\]
	and thus we choose a primitive 5th root of unity (say $\xi =\exp (2\pi i/5)$). Then the solutions are
\[(\xi^{1} ,\xi ^2,\xi ^3,\xi^{4})\]
We see that this group is homomorphic to ${\bm Z}/5{\bm Z}$, and that there's a certain relationship between the roots of this polynomial and a certain group.
\end{rmk}
\vspace{2ex}
\section{Chapter 13. Field Extensions}
\begin{defi}
(13.1) Let $K$ be a field. A field extension of $K$ is another field $F$ endowed with an injective ring homomorphism $\phi :K\rightarrow F$. You are essentially identifying the field with a subfield of a larger ring $F$.
\end{defi}
\vspace{2ex}
\begin{rmk}
From now on, we use the following conventions:
\begin{itemize}
	\item[(i)] Every ring is a commutative ring with $1$
	\item[(ii)] A ring homomorphism sends $1$ to $1$ 
\end{itemize}
\end{rmk}
\vspace{2ex}
\begin{rmk}
There is at most one ring homomorphism between two fields. Notice that $\mathop{\mathrm{ker}}\phi $ is an ideal (additive subgroup with the absorption property) of $K$, so that $\mathop{\mathrm{ker}}\phi $ is either $0$ or $K$. Therefore, $\phi $ must be injective.
\end{rmk}
\vspace{2ex}
\begin{itemize}

\end{itemize}

